\section{Descrição}\label{sec:descricao}

Nesta seção devem ser descritos os \textbf{aspectos mais relevantes de
especificação e implementação} para a compreensão sobre o trabalho
desenvolvido. O título ``DESCRIÇÃO'' pode ser complementado com ``DO
SOFTWARE'', ``DA FERRAMENTA'' ou ``DO PROTÓTIPO'' ou aquilo que melhor
representar o que foi desenvolvido. Esta seção deve estar organizada em
pelo menos duas subseções: especificação e implementação

Reitera-se que, em função da limitação do número de páginas, a descrição
deve contemplar o que é mais significativo para a compreensão do que foi
desenvolvido.

\subsection{Especificação}\label{subsec:especificacao}

\textbf{\ul{Os diagramas inseridos devem considerar o modelo de
estrutura e de comportamento do que foi desenvolvido}}. Se é um sistema
com usuário final, inclua a lista de requisitos funcionais e não
funcionais. Destaca-se que os \ul{diagramas desenvolvidos bem como
outros aspectos de especificação deverão obrigatoriamente constar nos
apêndices \textbf{quando não couberem nesta seção}}.

\subsection{Implementação}\label{subsec:implementacao}

Descreve os aspectos fundamentas para a compreensão da implementação do
TCC. Considere a inserção dos códigos mais relevantes da implementação,
bem como as telas do trabalho desenvolvido.
